\section{The endpoint failure and its logical scope}\label{sec:conclusion}

\begin{proof}[Proof of Theorem~\ref{thm:main}]
Work in ZFC with CH. Theorem~\ref{aec:class} gives an AEC \(K\) with
\(\LS(K)=\aleph_0\). Hence
\[
 H(K)=\beth_{(2^{\aleph_0})^+}=\beth_{\omega_2}
      =\beth_\eta.
\]
Theorem~\ref{end:two-models} supplies two nonisomorphic models of this
cardinality. Theorem~\ref{cat:tail} gives categoricity in every
\(\mu\ge\Lambda\). Finally, Cantor's theorem gives
\(\xi=2^\delta\ge\delta^+=\eta\), so \(\xi^+>\eta\). By strict
increase of the beth hierarchy,
\[
 \Lambda=\beth_{\xi^+}>\beth_\eta=H(K).
\]
Thus \(\lambda=\Lambda^+\) satisfies the premise of the proposed
transfer and its conclusion fails at \(\mu=H(K)\), as claimed.
\end{proof}

\subsection{The set-theoretic formulation}

For a literal first-order consistency statement, let \(\Theta\) assert
that there is a set
\(p=(\iota_k:k\ge1)\) of coding maps as in Lemma~\ref{col:coding}
for which the particular formulas of Definition~\ref{aec:objects}
define an AEC \(K_p\) with \(\LS(K_p)=\aleph_0\), two
nonisomorphic models of cardinality \(h(\aleph_0)\), and
categoricity in \(\Lambda^+>h(\aleph_0)\), where
\(\Lambda=\beth_{(2^{\aleph_1})^+}\).
Here the fixed definitions and the AEC and categoricity conditions
are expanded into one sentence of set theory; all quantified
parameters, structures, maps, and sequences are sets.

For the statements involving \(\Phi\), fix a first-order formulation
of the universal assertion~\eqref{eq:target} such that ZFC proves
its implication to the following fixed-family instance: for every
valid \(p\), if \(K_p\) is an AEC with \(\LS(K_p)=\aleph_0\),
\(\Lambda^+>h(\aleph_0)\), and categoricity in \(\Lambda^+\),
then \(K_p\) is categorical in \(h(\aleph_0)\).
The conclusions about \(\Phi\) below are read separately for each
such formulation. In particular, ZFC proves
\(\Theta\Rightarrow\neg\Phi\).

\begin{corollary}\label{cor:consistency}
ZFC proves \(\mathrm{CH}\Rightarrow\Theta\).
For each fixed \(\Phi\) as above, it also proves
\(\mathrm{CH}\Rightarrow\neg\Phi\), and
\[
 \Con(\mathrm{ZFC})
 \quad\Longrightarrow\quad
 \Con(\mathrm{ZFC}+\Theta)
 \quad\Longrightarrow\quad
 \Con(\mathrm{ZFC}+\neg\Phi).
\]
If ZFC is consistent, no such \(\Phi\) is a theorem of ZFC.
\end{corollary}

\begin{proof}[Proof of Corollary~\ref{cor:consistency}]
The class and strong relation are given by fixed formulas with the
set of coding maps as parameter. Expanding these formulas in the
AEC axioms and the stated model assertions gives the sentence
\(\Theta\). The construction and its verification in
Theorem~\ref{thm:main} therefore prove, in ZFC,
\(\mathrm{CH}\Rightarrow\Theta\).

For each \(\Phi\) specified above, its ZFC-provable fixed-family
instance applies to the parameter witnessing \(\Theta\).
It would give categoricity in \(h(\aleph_0)\), contradicting
the two nonisomorphic models required by \(\Theta\).
Thus ZFC proves \(\Theta\Rightarrow\neg\Phi\).

G\"odel's relative-consistency theorem
\cite{Godel1938} gives
\(\Con(\mathrm{ZFC})\Rightarrow
  \Con(\mathrm{ZFC}+\mathrm{CH})\).
Since the latter theory proves \(\Theta\), its consistency implies
\(\Con(\mathrm{ZFC}+\Theta)\), and the proved implication
\(\Theta\Rightarrow\neg\Phi\) then gives
\(\Con(\mathrm{ZFC}+\neg\Phi)\).
If \(\Phi\) were a theorem of consistent ZFC, adjoining
\(\neg\Phi\) could not give a consistent theory.
\end{proof}

The corollary gives one-sided nonprovability, not independence in both
directions: neither an unconditional ZFC proof of \(\neg\Phi\) nor
consistency of \(\Phi\) is asserted. The constructed class is categorical
on a higher tail, so the obstruction concerns the prescribed endpoint
and is compatible with qualitative eventual categoricity. No qualitative
eventual-categoricity theorem is used in the argument.

\subsection{Why structural transfer theorems do not apply}

For clarity, the construction fails both joint embedding and
amalgamation, although neither failure was needed as an input to the
proof. Joint embedding means that any two models admit strong
embeddings into a common model. Amalgamation requires such embeddings
over a given common strong submodel, agreeing on that submodel.

\begin{proposition}\label{prop:no-amalgamation}
The class \(K\) has neither joint embedding nor amalgamation.
The failure of amalgamation can be witnessed over a countable model.
\end{proposition}

\begin{proof}
Let \(M\) be the model with uncountable base from
Theorem~\ref{end:two-models}. By Theorem~\ref{aec:class}, choose a
countable \(A\le_K M\). Construct an extension \(N\) of \(A\) by
enlarging only its index sort to cardinality \(\Lambda^+\), keeping
\(B^N=B^A\), both old group sorts, and all fibers and bits belonging
to old tuples. For each new tuple add the required independent torsors
and any bit array satisfying the translation law. The countability of
\(B^N\) makes the object condition automatic. This is a basic
inclusion \(A\subseteq N\), and the strong condition is vacuous
because the base order has not grown. Thus \(A\le_K N\).

If \(M\) and \(N\) strongly embedded into a common model \(P\),
the image of \(M\) would force \(B^P\) to be uncountable and the image
of \(N\) would force \(|I^P|\ge\Lambda^+\).
This contradicts Proposition~\ref{cat:bounded}.
There is therefore no joint embedding of \(M,N\), and in particular
no amalgam of the two extensions \(A\le_K M,N\).
\end{proof}

Proposition~\ref{prop:no-amalgamation} explains why results imposing
amalgamation, such as \cite[Corollary~9.10]{Vasey2017Downward}, do not
contradict Theorem~\ref{thm:main}. Likewise, the existence Hanf theorem
does not imply that the particular model \(M\) has arbitrarily large
extensions. The full class has arbitrarily large models with countable
base; its uncountable-base part has the separate bound proved in
Proposition~\ref{cat:bounded}. That part cannot itself have
L\"owenheim--Skolem number \(\aleph_0\), because it has no countable
models.
